% ------------------------------------------------------------------------
% Relatório Parcial - PIBIT 2025 (versão revisada)
% ------------------------------------------------------------------------
\documentclass[
	12pt,
	openright,
	oneside,
	a4paper,
	brazil
]{abntex2}

% ----------------------------------------------------------
% PACOTES
% ----------------------------------------------------------
\usepackage{lmodern}
\usepackage[T1]{fontenc}
\usepackage[utf8]{inputenc}
\usepackage{indentfirst}
\usepackage{color}
\usepackage{graphicx}
\usepackage{microtype}
\usepackage{fancyhdr}
\usepackage[brazilian,hyperpageref]{backref}
\usepackage[alf]{abntex2cite}
\usepackage{hyperref}
\usepackage{multirow}
\usepackage{booktabs}
\usepackage{longtable}
\usepackage{array}

% ----------------------------------------------------------
% CONFIGURAÇÕES DE PACOTES
% ----------------------------------------------------------
\renewcommand{\backrefpagesname}{Citado na(s) página(s):~}
\renewcommand{\backref}{}
\renewcommand*{\backrefalt}[4]{
	\ifcase #1 Nenhuma citação no texto.%
	\or Citado na página #2.%
	\else Citado #1 vezes nas páginas #2.%
	\fi}

% ---
% Informações para CAPA e FOLHA DE ROSTO
% ---
\titulo{Relatório Parcial\\Desenvolvimento de um Dashboard para o Monitoramento e Acompanhamento Contínuo de Projetos em Fábricas de Software Acadêmicas}
\autor{Danrley Willyan da Silva Pereira (RGM 42570166)}
\orientador{Prof.ª Kerlla de Souza Luz}
\local{Brasília}
\data{2025}
\instituicao{
  Centro Universitário do Distrito Federal (UDF)\\
  Programa Institucional de Bolsas de Iniciação em Desenvolvimento Tecnológico e Inovação (PIBIT)
}

% ---
% Configurações de hyperref
% ---
\makeatletter
\hypersetup{
	pdftitle={\@title}, 
	pdfauthor={\@author},
	pdfsubject={Relatório Parcial PIBIT},
	pdfkeywords={Dashboard, Fábrica de Software Acadêmica, Métricas, KPIs, Arquitetura Medalhão},
	colorlinks=true,
	linkcolor=blue,
	citecolor=blue,
	filecolor=magenta,
	urlcolor=blue,
}
\makeatother

% Espaçamento
\setlength{\parindent}{1.3cm}
\setlength{\parskip}{0.2cm}

\makeindex

% ----
% Início do documento
% ----
\begin{document}
\selectlanguage{brazil}
\frenchspacing

% ----------------------------------------------------------
% ELEMENTOS TEXTUAIS
% ----------------------------------------------------------
\textual

% ===============================
% RELATÓRIO PARCIAL (PIBIT/2025)
% ===============================

\section*{1.\ Identificação do Projeto}
\textbf{Professor(a) Proponente:} Kerlla de Souza Luz \\
\textbf{Aluno(a):} Danrley Willyan da Silva Pereira (RGM 42570166) \\
\textbf{Área de conhecimento:} Engenharia de Software / Engenharia de Dados \\
\textbf{Título:} Desenvolvimento de um Dashboard para o Monitoramento e Acompanhamento Contínuo de Projetos em Fábricas de Software Acadêmicas

\section*{2.\ Resumo (máx.\ 2000 caracteres)}
Este relatório parcial apresenta os avanços do projeto de um \textit{dashboard} open source para monitorar projetos em Fábricas de Software Acadêmicas. Foram consolidados fichamentos da literatura e um conjunto inicial de \textbf{KPIs} técnicos e colaborativos; implementou-se uma \textbf{pipeline de dados} baseada na \textbf{Arquitetura Medalhão} com \textbf{ingestão diária} (Bronze), \textbf{modelagem Silver} e \textbf{publicação semanal} de indicadores (Gold). Destaca-se que \textbf{a adoção da Arquitetura Medalhão e dos ciclos de iteração com gestores/professores não constavam do escopo original} e foram incorporados após entrevistas qualitativas (Apêndices A e B), que também sustentaram a decisão de \textbf{não operar mais em ``tempo real''}. Em lugar disso, a atualização semanal do painel atende às necessidades pedagógicas e de governança. Os \textbf{testes-piloto} com equipes \textbf{ainda não foram iniciados}; antes, será conduzido um ciclo iterativo de publicação de dados/visões e coleta de \textit{feedback} para priorização dos próximos KPIs e visualizações.

\smallskip
{\footnotesize Estrutura do relatório conforme roteiro COP/UDF.}

\section*{3.\ Mudanças de escopo e impactos}
\textbf{(M1) Arquitetura de dados.} A proposta original não previa o uso da \textbf{Arquitetura Medalhão}. Foi incorporada para garantir \textit{linhagem}, reprodutibilidade e abertura à colaboração: \textbf{Bronze} (ingestão diária, \emph{append-only}); \textbf{Silver} (normalização em \texttt{dim}/\texttt{fact}); \textbf{Gold} (KPIs semanais para o painel). \\
\textbf{Impacto:} maior qualidade e auditabilidade dos dados; possibilidade de reutilização por outras fábricas (licença copyleft). 

\textbf{(M2) Periodicidade.} A proposta mencionava \emph{tempo real/quase real}. Após as entrevistas, \textbf{descontinuou-se a coleta em tempo real}; \textbf{ingestão diária} e \textbf{publicação semanal} mostraram-se suficientes. \\
\textbf{Impacto:} redução de custo/complexidade, sem prejuízo pedagógico.

\textbf{(M3) Iterações com stakeholders.} A proposta não previa \textbf{ciclos de co-desenho} com gestores e docentes. Passamos a operar \textbf{iterações curtas}: publicar dados/visões, colher \textit{feedback}, priorizar próximos KPIs. \\
\textbf{Impacto:} maior aderência às necessidades do semestre; expectativas realistas para avaliação.

\textbf{(M4) Testes-piloto.} \textbf{Ainda não iniciados}. Serão substituídos, no curto prazo, por \textbf{rodadas iterativas} (dados/visões + \textit{feedback}) e, em seguida, estruturados como pilotos formais.

\section*{4.\ Objetivos (status parcial)}
{\footnotesize Legenda: \textbf{T} concluído; \textbf{P} em progresso; \textbf{N} não iniciado; \textbf{$\Delta$} escopo ajustado.}
\begin{enumerate}
    \item \textbf{Revisão de literatura (OE1)}: \textbf{P} \,\,(\em fichamentos consolidados; seção~\ref{sec:fichamentos}). 
    \item \textbf{Definir KPIs (OE2)}: \textbf{P}$\Delta$ \,\,(\em conjunto inicial definido; evolução passará a seguir ciclos de \textit{feedback}). 
    \item \textbf{Protótipo de coleta (OE3)}: \textbf{P}$\Delta$ \,\,(\em ETL na Arquitetura Medalhão adotada; \textbf{ingestão diária}). 
    \item \textbf{Construir o dashboard (OE4)}: \textbf{P}$\Delta$ \,\,(\em publicação \textbf{semanal} de KPIs; co-desenho com gestores/docentes). 
    \item \textbf{Testes-piloto (OE5)}: \textbf{N} \,\,(\em serão iniciados após 1--2 iterações de \textit{feedback} com usuários).
\end{enumerate}

\section*{5.\ Resultados obtidos até o momento}
\textbf{Pipeline de dados.} 
\emph{Bronze}: coleta diária via API do GitHub (organização, membros, repositórios, commits, issues, PRs, reviews/comentários; quando disponível, \emph{traffic/clones}); filtros de qualidade (remoção de \emph{forks} e \emph{blacklists}); \emph{append-only}. \\
\emph{Silver}: normalização de JSONs em \texttt{dim}/\texttt{fact} (\texttt{dim\_users}, \texttt{dim\_repos}, \texttt{fact\_commits}, \texttt{fact\_issues}, \texttt{fact\_prs}, \texttt{fact\_reviews}); chaves consistentes; padronização temporal semanal; colunas de linhagem/auditoria; relações autor$\leftrightarrow$revisor. \\
\emph{Gold}: KPIs iniciais publicados semanalmente (atividade de código, fluxo de mudanças, revisões, colaboração, regularidade, concentração), incluindo \textbf{Slope\_8w}, \textbf{Clones\_W1--W2}, \textbf{centralidade} em grafos e \emph{cycle time}.

\textbf{Visualizações exploratórias.} 20 gráficos úteis para leitura de padrões de engajamento, colaboração e fluxo (lista no \textbf{Apêndice C}). 

\textbf{Achados provisórios.} \emph{Ramp-up} de engajamento na primeira metade do semestre; \emph{clones} iniciais como sinal de adesão; distinção de perfis \emph{integrador} vs.\ \emph{especialista} por grafos; suficiência de atualização semanal para o uso pedagógico.

\section*{6.\ Análises e discussões}
As evidências qualitativas (Apêndices A e B) contextualizam as séries: (i) baixas iniciais de atividade não devem ser penalizadas antes do \emph{ramp-up}; (ii) \emph{clones} iniciais ajudam a estimar adesão; (iii) grafos autor$\leftrightarrow$revisor revelam papéis. A publicação \textbf{semanal} (Gold) dá previsibilidade para docentes/gestores, enquanto \textbf{Bronze/Silver} permanecem abertos a análises avançadas por usuários experientes.

\section*{7.\ Fichamentos da revisão bibliográfica e impacto nos KPIs}
\label{sec:fichamentos}
Foram produzidos \textbf{fichamentos} das referências centrais do projeto, registrando: objetivo do trabalho, método, principais métricas/indicadores e limitações. A Tabela~\ref{tab:fichamentos-kpi} resume como esses fichamentos ancoram o desenho dos KPIs (OE2).

\begin{table}[h]
\centering
\caption{Relação entre fichamentos (OE1) e KPIs adotados (OE2)}
\label{tab:fichamentos-kpi}
\renewcommand{\arraystretch}{1.15}
\begin{tabular}{p{5cm} p{9cm}}
\toprule
\textbf{Fichamento (referência)} & \textbf{Direcionamento para KPIs} \\
\midrule
Kitchenham (2010) & Reforça planos de medição claros (objetivo, fonte, periodicidade, responsáveis) e triangulação; sustenta \emph{regularidade semanal}, transparência/linhagem e cautela com leitura isolada de contagens. \\
Hamer et al.\ (2021) & Apoia atividade de código, \emph{code churn}, PRs, desigualdade/concentração e análise de colaboração (grafos). \\
Barbosa \& Brum (2016) & Orienta indicadores gerenciais: throughput de \emph{issues}, tempo de resolução, aderência de processo (atribuições/estados). \\
Oliveira et al.\ (2017); dos Santos et al.\ (2021) & Contextualizam métrica em fábricas acadêmicas; importância de feedback formativo e integração prática-teoria. \\
Ozkan \& Mishra (2019); Temitope (2020) & Enfatizam visualização e responsividade (tempo até 1ª resposta), integração de ferramentas e monitoramento contínuo. \\
\bottomrule
\end{tabular}
\end{table}

\subsection{Quadro-s{\'\i}ntese da literatura e rela{\c c}{\~a}o com a arquitetura medalh{\~a}o}
A Tabela~\ref{tab:quadro-sintese-medalhao} organiza as principais fontes, seus objetivos e contribui{\c c}{\~o}es para o desenho de KPIs e do pipeline em camadas (Bronze~$\rightarrow$~Silver~$\rightarrow$~Gold).

\begin{longtable}{p{3.5cm} p{4.5cm} p{6cm}}
\caption{Síntese da literatura: contribuições para KPIs e arquitetura} \\
\label{tab:quadro-sintese-medalhao} \\
\toprule
\textbf{Fonte} & \textbf{Objetivo/Contexto} & \textbf{Contribuições para KPIs/Processo} \\
\midrule
\endfirsthead

\caption[]{(continuação)} \\
\toprule
\textbf{Fonte} & \textbf{Objetivo/Contexto} & \textbf{Contribuições para KPIs/Processo} \\
\midrule
\endhead

\cite{barbosaIndicadoresConjunto2016} & Guia de plano de medição e catálogo de indicadores & Base para dicionário de KPIs e governança da medição (Silver→Gold) \\
\midrule
\cite{barbosaGeracaoGitHub} & Geração de informações gerenciais do Git/GitHub & ETL e views semânticas para KPIs de atividade/fluxo (Bronze→Silver→Gold) \\
\midrule
\cite{oliveira2022plataforma} & Plataforma de coleta/análise do GitHub & Facilita ingestão e persistência; alimenta normalização (Bronze→Silver) \\
\midrule
\cite{vega2022scrumwatch} & Monitorar equipes Scrum com integração multi-fonte & Métricas de sprint e qualidade; consultas sobre fatos/dimensões (Silver→Gold) \\
\midrule
\cite{hamer2021using} & Métricas Git em projetos estudantis & KPIs de atividade, churn, distribuição e centralidade (Silver→Gold) \\
\midrule
\cite{kitchenham2010s} & Princípios de métricas e medição & Diretrizes de validade/periodicidade; regularidade como dimensão (Gold) \\
\midrule
\cite{ozkan2019agile} & Ferramentas de gestão ágil & KPIs de fluxo (PR/issue), revisões e aderência a processo (Gold) \\
\midrule
\cite{temitope2020software} & Adoção de software de gestão & Impactos em eficiência e colaboração; métricas de revisões/tempo (Gold) \\
\midrule
Fábricas Acadêmicas\textsuperscript{*} & Experiências e modelos pedagógicos & Demanda por visibilidade/indicadores; leitura pedagógica de KPIs \\
\bottomrule
\multicolumn{3}{l}{\footnotesize \textsuperscript{*} \cite{dos2021experiencia,romanha2016modelo,romanha2015fabrica}} \\
\end{longtable}


\noindent \textbf{Resultado:} o conjunto inicial de KPIs inclui: atividade (\emph{commits/semana}, PRs), fluxo (\emph{lead time} de PR, tempo até 1ª resposta), backlog (issues abertas/fechadas, \emph{aging}), colaboração (participação em \emph{reviews}, grafos), regularidade (semanas ativas) e concentração (desigualdade). Indicadores motivados pelos relatos (\textbf{Slope\_8w}, \textbf{Clones\_W1--W2}, \textbf{centralidade}) foram \emph{adicionados} por pertinência pedagógica.

\section*{8.\ Próximas atividades (próxima iteração)}
\textbf{Dados}: ampliar cobertura Bronze e ajustes de chaves/dicionário Silver. \\
\textbf{KPIs e painel}: incorporar métricas priorizadas em \textit{feedback} (Gold), com páginas por time/repo/indivíduo e módulo de grafos. \\
\textbf{Stakeholders}: conduzir 1--2 ciclos de \textit{feedback} com gestores e docentes. \\
\textbf{Pilotos}: iniciar \textbf{testes-piloto} após as iterações iniciais, com protocolo de avaliação (usabilidade, utilidade, aderência).

\section*{9.\ Cronograma atualizado (alto nível)}
\begin{itemize}
    \item \textbf{Set–Out/2025}: ingestão diária estável (Bronze), ajustes Silver, painel v0.2, 1ª rodada de \textit{feedback}.
    \item \textbf{Nov/2025}: KPIs priorizados (Gold), módulo de grafos, painel v0.3, 2ª rodada de \textit{feedback}.
    \item \textbf{Dez/2025}: consolidação de análises e preparação dos \textbf{testes-piloto} (jan/2026).
\end{itemize}

% ----------------------------------------------------------
% REFERÊNCIAS BIBLIOGRÁFICAS
% ----------------------------------------------------------
\bibliography{relatorio-parcial}

% ----------------------------------------------------------
% APÊNDICES
% ----------------------------------------------------------
\begin{apendicesenv}
\partapendices

\chapter{Entrevista com LAPPIS/UnB (FGA)}
\label{apendiceA}

\textbf{Participantes:} Danrley Pereira (IBICT) e Dra.\ Carla (Coordenadora LAPPIS) \\
\textbf{Data:} 22/08/2025 \\
\textbf{Local:} FGA -- UnB \\
\textbf{Tema:} Fábricas de software acadêmicas, métricas de interação e experiências do LAPPIS

\section*{1.\ Contexto Institucional}
\begin{itemize}
\item Envolvimento atual em \textbf{sistemas estruturantes} para controle orçamentário no poder executivo.
\item Desafio de \textbf{fragmentação de sistemas} e baixa interoperabilidade.
\end{itemize}

\section*{2.\ Reflexões sobre o modelo de ``Fábrica de Software''}
\begin{itemize}
\item Crítica ao conceito tradicional para ambientes acadêmicos/inovadores; necessidade de \textbf{repensar e adaptar} práticas ao aprendizado.
\end{itemize}

\section*{3.\ Arquitetura e tratamento de dados}
\begin{itemize}
\item Proposta da \textbf{Arquitetura Medalhão}:
\begin{itemize}
\item \textbf{Bronze (Raw):} extração de GitHub, GitLab, Notion;
\item \textbf{Silver:} transformação e homogeneização em modelos definidos;
\item \textbf{Gold:} camada de conhecimento com insights consolidados.
\end{itemize}
\item \textbf{Atualização semanal} considerada suficiente para o dashboard.
\end{itemize}

\section*{4.\ Métricas Quantitativas e Qualitativas}
\begin{itemize}
\item Necessidade de \textbf{interpretar quantitativos via análises qualitativas}.
\end{itemize}

\section*{5.\ Formação Acadêmica e Experiência}
\begin{itemize}
\item Disciplinas regulares: \textbf{4 meses}; alunos ingressam a partir do \textbf{4º semestre} e permanecem, em média, \textbf{2 anos}.
\end{itemize}

\section*{6.\ Métricas de Interações em Repositórios}
\begin{itemize}
\item Potencial de \textbf{grafos de rede} para mapear perfis:
\begin{itemize}
\item \textit{Aluno integrador}: muitos vínculos entre membros/repos;
\item \textit{Aluno especialista}: foco em áreas específicas (ex.: SRE/DevOps).
\end{itemize}
\end{itemize}

\section*{7.\ Conclusões e Próximos Passos}
\begin{itemize}
\item Alinhamento sobre utilidade de \textbf{métricas de interação} e adoção do modelo em \textbf{camadas (Bronze~$\rightarrow$~Silver~$\rightarrow$~Gold)}.
\item Reforço da \textbf{integração de contexto qualitativo} às análises.
\end{itemize}

\chapter{Relatos LABTECH/UDF}
\label{apendiceB}

\section*{Resumo dos relatos}
\begin{itemize}
\item No início do semestre, \textbf{interações no GitHub tendem a ser tímidas}, pois os alunos ainda estão se ambientando ao projeto e ao domínio.
\item A \textbf{curva inicial} pode ser mais baixa, com \textbf{aceleração ao redor da metade do semestre}.
\item \textbf{Commits e interações} mostram quem efetivamente contribuiu e buscou engajar.
\item O \textbf{número de \textit{clones}} no começo do semestre sinaliza a \textbf{intenção de participação}; acompanhar sua evolução pode indicar adesão.
\end{itemize}

\section*{Implicações para os KPIs}
\begin{itemize}
\item Monitorar a \textbf{inclinação inicial} (semanas 1--8) como possível preditor de sucesso posterior.
\item Persistir e analisar \textbf{clones} no início (\emph{traffic/clones}) como indicador de intenção.
\item Separar leituras por \textbf{papéis} (autor, revisor) e por \textbf{áreas} (dados, frontend, APIs, plataforma), permitindo distinguir \textit{integradores} de \textit{especialistas}.
\end{itemize}


\chapter{Análise de Dados GitHub do LabTech}
\label{apendiceC}

\section*{Contexto}
Este apêndice resume a exploração do diretório \texttt{/assets} e a execução do script \texttt{extrair-dados-github.py} para coleta e análise de dados da organização \texttt{LabTechUDF} no GitHub, servindo como base empírica do artigo.

\section*{Como os dados foram capturados}
A pipeline implementa um \textbf{ETL} completo com:
\begin{itemize}
    \item Coleta via \textbf{GitHub API} da organização \texttt{LabTechUDF};
    \item \textbf{Cache} para evitar chamadas desnecessárias e \textbf{rate limit monitoring};
    \item \textbf{Filtragem} automática (remoção de \textit{forks} e \textit{blacklist} de repositórios);
    \item Coletas abrangentes: dados organizacionais e membros; repositórios e colaboradores; \textit{issues}, \textit{pull requests} e eventos; \textit{commits} e \textit{timeline}; \textbf{mapeamento de relacionamentos} entre contribuidores.
\end{itemize}

\section*{Arquitetura de dados (Medalhão)}
\begin{itemize}
    \item \textbf{Bronze (raw)}: ingestão \emph{append-only} dos JSONs de origem (GitHub), preservando \textit{payloads}, carimbos de tempo e metadados;
    \item \textbf{Silver (cleaned/modelado)}: normalização em modelos tabulares (\texttt{dim\_repos}, \texttt{dim\_users}, \texttt{fact\_commits}, \texttt{fact\_prs}, \texttt{fact\_issues}, \texttt{fact\_reviews}), com padronização temporal e colunas de linhagem;
    \item \textbf{Gold (métricas/KPIs)}: agregações para o \textit{dashboard} (throughput de \textit{issues}, \textit{lead time} de PR, \textit{code churn} semanal, participação em \textit{reviews}).
\end{itemize}
\noindent Atualização \textbf{semanal}, suficiente para fins pedagógicos, mantendo \textbf{abertura à colaboração}: usuários avançados podem consumir Bronze/Silver para análises multi-relacionadas, enquanto o Gold provê KPIs prontos.

\section*{Métricas extraídas}
\textbf{Técnicas:} commits por usuário/repositório/tempo; \textit{issues} criadas/atribuídas/fechadas; \textit{PRs} abertos/revisados; \textit{cycle time} para \textit{issues}/\textit{PRs}; \textit{burndown}; \textit{code churn} e frequência de commits.

\noindent\textbf{Colaborativas:} rede de colaboração; \textit{cross-collaboration} (múltiplos repositórios); participação em \textit{reviews}/comentários; centralidade e papéis (integrador vs.\ especialista).

\noindent\textbf{Comportamentais:} \textit{activity heatmaps} (dia x hora); \textit{maturity score} (idade de conta, repositórios públicos, seguidores); perfil de membro (novo vs.\ estabelecido).

\section*{Visualizações geradas (20)}
\begin{enumerate}
    \item \textit{activity-heatmap.png} — Heatmap de atividade por dia/hora
    \item \textit{burndown-chart.png} — Burndown de \textit{issues}
    \item \textit{collaboration-network.png} — Rede de colaboração completa
    \item \textit{collaboration-network-no-isolated.png} — Rede sem nós isolados
    \item \textit{commit-boxplot.png} — Boxplot de distribuição de commits
    \item \textit{commit-frequency.png} — Frequência de commits
    \item \textit{contributions-over-time.png} — Contribuições ao longo do tempo
    \item \textit{contributions-over-time-filtered.png} — Contribuições (Aug--Dec 2024)
    \item \textit{contributions-per-member.png} — Contribuições por membro
    \item \textit{contributions-per-repo.png} — Contribuições por repositório
    \item \textit{cross-collab-histogram.png} — Colaboração cruzada
    \item \textit{cycle-time-histogram.png} — Tempo de resolução de \textit{issues}
    \item \textit{established-contributions.png} — Membros estabelecidos vs.\ contribuições
    \item \textit{github-usage-heatmap.png} — Heatmap de uso do GitHub
    \item \textit{issue-states.png} — Estados de \textit{issues} (abertas vs.\ fechadas)
    \item \textit{maturity-vs-contributions.png} — Maturidade vs.\ contribuições
    \item \textit{new-pie-contribution.png} — Novos membros vs.\ contribuições
    \item \textit{pie-contributions.png} — Distribuição geral de contribuições
    \item \textit{pr-cycle-time-histogram.png} — Tempo de resolução de PRs
    \item \textit{repos-followers.png} — Repositórios públicos vs.\ seguidores
\end{enumerate}

\section*{Principais insights}
\begin{itemize}
    \item \textbf{Arquitetura medalhão} efetivamente implementada (Bronze~$\rightarrow$~Silver~$\rightarrow$~Gold);
    \item \textbf{Triangulação qualitativa} com os \textbf{Apêndices A e B} para interpretar variações das séries;
    \item KPIs específicos como \textbf{Slope\_8w} (ritmo inicial) e \textbf{centralidade} em grafos de interação;
    \item \textbf{Atualização semanal} suficiente para governança pedagógica;
    \item Identificação de \textbf{hubs} (membros que contribuem em múltiplos repositórios) e padrões de \textit{cross-collaboration}.
\end{itemize}

\section*{Reprodutibilidade e ética}
\begin{itemize}
    \item Artefatos em \texttt{/assets} e \texttt{extrair-dados-github.py}; logs de \textit{rate limit} e \textit{cache};
    \item Linhagem preservada (colunas de auditoria no Silver); KPIs recomputáveis;
    \item Anonimização de identificadores para divulgação acadêmica, conforme diretrizes éticas do projeto.
\end{itemize}


\end{apendicesenv}

\end{document}
